% Generated draft: requires mathematical and source review.
\documentclass[11pt]{article}
\usepackage[margin=25mm]{geometry}
\usepackage{fontspec}
\setmainfont{Latin Modern Roman}
\usepackage{amsmath,amssymb,mathtools,unicode-math}
\setmathfont{Latin Modern Math}
\usepackage{enumitem,needspace,xurl,hyperref}
\setlength{\parindent}{0pt}
\setlength{\parskip}{5pt}
\setlength{\emergencystretch}{3em}
\newcommand{\N}{\mathbb{N}}
\newcommand{\Z}{\mathbb{Z}}
\newcommand{\Q}{\mathbb{Q}}
\newcommand{\R}{\mathbb{R}}
\newcommand{\C}{\mathbb{C}}
\newcommand{\F}{\mathbb{F}}
\newcommand{\A}{\mathbb{A}}
\newcommand{\PP}{\mathbb{P}}
\newcommand{\E}{\mathbb{E}}
\newcommand{\eps}{\varepsilon}
\newcommand{\dd}{\,\mathrm{d}}
\begin{document}
\section*{UnsolvedMath ID: 20003107}
\textbf{Source problem code:} catalogue entry
\subsection{Definitions and mathematical statement}
K. Rafi asks whether the Thurston metric is a degeneration of a better vector-valued symmetric distance on Teichm\"uller space. The extracted record also defines the reversed metric $d^*(X,Y)=d_{\mathrm{Th}}(Y,X)$; trailing dots are layout. The words ``better,'' ``symmetric'' for a vector value, and ``degeneration'' are not specified. The aid distinguishes literal symmetry from symmetry involving an opposition involution.

Write $d=d_{\mathrm{Th}}$ in the finite-type marked hyperbolic setting and define
\[
 D(X,Y)=(d(X,Y),d(Y,X))\in\R_{\geq0}^2,
 \qquad\iota(a,b)=(b,a).
\]
Then $D(Y,X)=\iota D(X,Y)$ and the triangle inequality holds coordinatewise. The aid calls this an opposition-symmetric cone-valued metric; it is generally not literally symmetric under interchange of $X,Y$.

Its curve-resolved refinement uses the signed log-length cocycle
\[
 L_\gamma(X,Y)=\log\ell_\gamma(Y)-\log\ell_\gamma(X),\qquad
 L_\gamma(X,Z)=L_\gamma(X,Y)+L_\gamma(Y,Z).
\]
Set $L_\gamma^+=\max(L_\gamma,0)$ and $L_\gamma^-=\max(-L_\gamma,0)$. The pair of nonnegative curve-indexed functions has coordinatewise triangle inequalities and is exchanged under reversal. Taking their supremum norms gives $D$ in the usual nonnegative Thurston-distance convention. More generally, if an opposed lift $E$ has an involution $\iota_E$ and scalar readout $f$ recovering $d$, the reported factor map on its realized values is $z\mapsto(f(z),f(\iota_Ez))$. This is a formal recovery property, not an independently characterized optimal geometric lift.

The aid reports symmetric scalarizations: for $1\leq p<\infty$,
\[
 d_p(X,Y)=\bigl(d(X,Y)^p+d(Y,X)^p\bigr)^{1/p},
\]
and $d_\infty=\max\{d,d^*\}$ are genuine symmetric metrics. A precise biased family, in the elementary linear convention, is
\[
 d_\theta=\theta d+(1-\theta)d^*,\qquad0\leq\theta\leq1.
\]
Its center $\theta=1/2$ is symmetric and its endpoint $\theta=1$ is $d$. The intermediate members need not be symmetric; this is not a family consisting solely of symmetric scalar metrics that converges pointwise to an asymmetric one.

Sorting the two coordinates gives the literally symmetric chamber value
\[
 C(X,Y)=\bigl(\max\{d,d^*\},\min\{d,d^*\}\bigr).
\]
The reported triangle property is weak majorization,
\[
 C(X,Z)\preceq_w C(X,Y)+C(Y,Z),
\]
meaning the first/largest coordinate and the sum of both coordinates on the left are bounded by the corresponding quantities on the right. A coordinatewise triangle inequality for the sorted pair is not asserted.

Conversely, a pointwise limit on a fixed set of literally symmetric scalar distances remains symmetric. Thus such a limit cannot equal a genuinely asymmetric Thurston distance. The aid concludes that a degeneration retaining asymmetry must keep opposition or oriented limiting information, or change the comparison framework. These constructions and the pointwise no-go statement are partial interpretations of the AIM request; no unique meaning of a better geometric vector-valued distance or a full degeneration theory is supplied.
\subsection{Short English statement}
The forward/reverse pair has symmetry by swapping its coordinates and admits symmetric scalar metrics. Biasing the readout recovers Thurston distance at an endpoint, but that family is not entirely symmetric. The aid distinguishes this construction from an impossible pointwise asymmetric limit of literally symmetric scalar distances.
\subsection{Sources}
\begin{itemize}
\item Local UnsolvedMath registry, permanent ID 20003107, source code catalogue entry; \texttt{aim-topology-notes.json}, notes[194], zero-based index 194.
\item AIM, \emph{Lipschitz metric on Teichmueller space}, Problem 3.4, K. Rafi: \url{https://aimath.org/pastworkshops/lipschitzteichproblems.pdf}. The supplied extraction includes the reversed-metric definition. No original bibliography is supplied.
\item Supplied research summary on opposition lifts, scalarizations, weak majorization and the symmetry obstruction; no independent full theorem citations are provided.
\item Machine-generated research attempt 1: partial result, open at the research run, rigor labelled ``proved,'' low novelty confidence; registry partially solved. Not peer reviewed; independently verify before citation. Registry updated 14 August 2026.
\end{itemize}
\end{document}
